% !TEX root = ../proyect.tex

\chapter{Conclusiones y desarrollos futuros}\label{cap:conclusiones}

Este capítulo cierra la memoria del trabajo. Tras un proyecto extenso conviene levantar la vista del detalle técnico y revisar qué se planteó al inicio, qué ha quedado efectivamente construido y qué queda por hacer. La estructura del capítulo reproduce ese orden: se sintetiza el estado final del sistema con un cuadro de indicadores, se contrastan los objetivos enunciados en el Capítulo~\ref{cap:introduccion} con la evidencia recogida durante la validación, se recopilan las lecciones aprendidas a lo largo del desarrollo, se analiza la desviación de la planificación inicial, se enumeran las líneas de evolución previstas para una segunda iteración del sistema, se discuten las consideraciones éticas y se cierra con una reflexión personal sobre la experiencia.

\section{Linceus en números}\label{sec:linceus-numeros}

Antes de revisar el cumplimiento de objetivos conviene fijar el tamaño y la cobertura efectiva del sistema en el momento del cierre. La Tabla~\ref{tab:linceus-numeros} agrupa los indicadores principales en cuatro bloques (alcance funcional, validación, métricas no funcionales y proceso de desarrollo). Todas las cifras son trazables al artefacto de origen citado en la columna derecha.

\begin{table}[hbtp]
\centering
\small
\caption{Indicadores cuantitativos de Linceus Assistant al cierre del prototipo (26/04/2026).}\label{tab:linceus-numeros}
\begin{tabular}{|p{5.4cm}|p{2.0cm}|p{4.6cm}|}
\hline
\textbf{Indicador} & \textbf{Valor} & \textbf{Fuente} \\ \hline
\multicolumn{3}{|l|}{\textit{Alcance funcional}} \\ \hline
Intenciones NLU definidas (activas) & 36 (34) & \texttt{domain.yml} \\ \hline
Entidades NLU & 10 & \texttt{domain.yml} \\ \hline
\textit{Stories} de entrenamiento & 37 & \texttt{data/stories.yml} \\ \hline
Reglas activas & 20 & \texttt{data/rules.yml} \\ \hline
\textit{Custom actions} activas & 14 & \texttt{actions/} \\ \hline
Archivos NLU por dominio & 5 & \texttt{data/nlu/} \\ \hline
Titulaciones soportadas & 3 (GII-IS, IC, TI) & Sección~\ref{sec:req-actores} \\ \hline
\multicolumn{3}{|l|}{\textit{Validación funcional}} \\ \hline
Casos E2E (PASS / total) & 150 / 159 (94,3~\%) & Tabla~\ref{tab:e2e-categorias} \\ \hline
RAG \textit{baseline} (PASS / total) & 47 / 50 (94,0~\%) & Sección~\ref{subsec:rag-baseline} \\ \hline
RAG profundidad (OK / total) & 68 / 74 (91,9~\%) & Tabla~\ref{tab:rag-cobertura} \\ \hline
NLU + Policy (acciones correctas) & 158 / 158 (100~\%) & Tabla~\ref{tab:nlu-resultados} \\ \hline
Tráfico real piloto (sesiones PASS) & 58 / 59 (98,3~\%) & Tabla~\ref{tab:trafico-resultados} \\ \hline
Suite admin \textit{pytest} (PASS / total) & 24 / 24 (100~\%) & Tabla~\ref{tab:admin-resultados} \\ \hline
Turnos auditados en piloto & 201 / 203 & Sección~\ref{subsec:trafico-resultados} \\ \hline
\multicolumn{3}{|l|}{\textit{Métricas no funcionales}} \\ \hline
Latencia mediana (p50) & 3,4 s & Tabla~\ref{tab:latencia} \\ \hline
Latencia p95 & 10,5 s & Tabla~\ref{tab:latencia} \\ \hline
Coste medio por turno (Gemini) & 0,012 cts & Sección~\ref{subsec:no-funcional-coste} \\ \hline
Llamadas a Gemini medidas & 216 & Sección~\ref{sec:pruebas-no-funcional} \\ \hline
Tiempo de cumplimiento RNF-14 (primera consulta) & 28 s (mediana) & Tabla~\ref{tab:rnf} \\ \hline
\multicolumn{3}{|l|}{\textit{Proceso de desarrollo}} \\ \hline
\textit{Sprints} completados & 8 & Capítulo~\ref{cap:planificacion} \\ \hline
Decisiones técnicas documentadas & 72 (D-001 a D-072) & \texttt{docs/sprints/} \\ \hline
Esfuerzo estimado total & $\sim$300 h & Sección~\ref{sec:conc-desviacion} \\ \hline
\end{tabular}
\end{table}

Los porcentajes superan en todas las capas su umbral académico de referencia con un margen razonable y los indicadores de proceso (72 decisiones documentadas, 8 \textit{sprints} cerrados con incremento desplegable) avalan que la cifra de validación es producto de un desarrollo trazable y no de un ajuste \textit{ad hoc} en la fase final.

\section{Cumplimiento de objetivos}\label{sec:conc-objetivos}

El Capítulo~\ref{cap:introduccion} planteaba un objetivo general y un conjunto de objetivos técnicos y académicos. Resulta razonable revisar uno por uno hasta qué punto se han alcanzado, apoyándose en la evidencia recopilada en el Capítulo~\ref{cap:pruebas}.

\subsection{Objetivo general}\label{subsec:conc-obj-general}

El objetivo general consistía en diseñar y desarrollar un asistente conversacional que ofreciera al alumnado de la Universidad de Sevilla un punto de acceso centralizado, accesible y eficiente a la información académica y administrativa. Tras el cierre del prototipo, este objetivo puede considerarse alcanzado en su versión correspondiente a la fase 1: el sistema cubre las tres épicas funcionales previstas (asignaturas, profesorado y horarios), responde correctamente al 94,3\,\% de los casos de la suite de aceptación funcional y al 98,3\,\% de las sesiones generadas por usuarios reales durante el piloto. La pieza queda desplegada en infraestructura de DigitalOcean y accesible mediante navegador, sin necesidad de credenciales, lo que satisface la condición de accesibilidad. Los grupos a los que la propuesta atendía con mayor prioridad (alumnado de nuevo ingreso e internacional) figuran entre los principales usuarios del piloto, lo que respalda la utilidad práctica del sistema.

\subsection{Objetivos técnicos}\label{subsec:conc-obj-tecnicos}

La Tabla~\ref{tab:cumplimiento-tecnico} cruza cada objetivo técnico definido en el Capítulo~\ref{cap:introduccion} con su evidencia en el Capítulo~\ref{cap:pruebas}. Todos ellos se han cumplido en la fase 1, con dos matices: la interfaz web final no se ha desarrollado en React como contemplaba la propuesta inicial sino en JavaScript \textit{vanilla} sobre HTML estático, decisión justificada por el menor coste de mantenimiento del prototipo, y el cumplimiento del RGPD se materializa por la naturaleza anónima de las sesiones más que por un despliegue \textit{on-premise}, dado que la infraestructura final reside en DigitalOcean y la base de datos en Supabase.

\begin{table}[hbtp]
\centering
\small
\caption{Cruce entre objetivos técnicos y evidencia obtenida.}\label{tab:cumplimiento-tecnico}
\begin{tabular}{p{6cm}p{7cm}}
\hline
\textbf{Objetivo técnico} & \textbf{Evidencia} \\
\hline
Sistema conversacional Rasa con flujos contextuales & Modelo Rasa 3.6 entrenado, 158/158 acciones correctas en la prueba de política (\S\ref{sec:pruebas-nlu}). \\
Integración con Gemini & 216 llamadas medidas en la suite, coste medio de 0,012 cts/turno (\S\ref{subsec:no-funcional-coste}). \\
Base de datos híbrida en Supabase con \texttt{pgvector} & Esquema relacional + tabla de \textit{embeddings}, capa RAG validada al 91,9\,\% de profundidad (\S\ref{subsec:rag-profundidad}). \\
Acciones personalizadas en Python & Tres épicas implementadas; 150/159 PASS en la suite E2E (\S\ref{sec:pruebas-e2e}). \\
\textit{Frontend} accesible & \textit{Widget} desplegado, validado mediante 59 sesiones reales (\S\ref{sec:pruebas-trafico}). \\
Cumplimiento RGPD & Sesiones anónimas, sin almacenamiento de datos personales identificables. \\
\hline
\end{tabular}
\end{table}

\subsection{Objetivos académicos y personales}\label{subsec:conc-obj-academicos}

Desde el punto de vista formativo, el proyecto ha permitido aplicar de forma articulada conocimientos procedentes de varias asignaturas del Grado en Ingeniería del Software (bases de datos, ingeniería del software, sistemas inteligentes, procesamiento del lenguaje natural y despliegue de aplicaciones), y ha exigido autoaprendizaje sobre tecnologías que no formaban parte del currículo oficial, como el módulo RAG y la integración de modelos de lenguaje. La metodología SCRUM adaptada a un proyecto individual ha sido también una novedad relevante, dado que el grado expone fundamentalmente a equipos de cinco o seis personas. El trabajo ha quedado documentado, validado mediante una suite de pruebas reproducible y desplegado en producción, lo que cumple los tres criterios formativos planteados al inicio: aplicar, integrar y dejar trazabilidad del proceso.

\section{Lecciones aprendidas}\label{sec:conc-lecciones}

La experiencia acumulada durante el desarrollo deja un conjunto de lecciones que resulta razonable consignar. Las primeras son de naturaleza técnica y se han ido formulando como decisiones documentadas a lo largo del repositorio (entradas \texttt{D-XXX}). Las segundas son de naturaleza metodológica y afectan al modo en que el proyecto se ha organizado en el tiempo. Las terceras son de naturaleza personal.

En el plano técnico, la lección más relevante es que la combinación de \textit{Text-to-SQL} y recuperación aumentada por generación cubre dominios distintos y conviene mantenerlas como caminos separados. \textit{Text-to-SQL} resuelve con eficacia las consultas estructuradas (créditos, listados, conteos), mientras que el RAG es indispensable cuando la pregunta requiere fragmentos textuales no normalizados (criterios de evaluación, composición de tribunales, contenidos de bloques). El intento inicial de unificar ambas vías mediante un único \textit{prompt} produjo respuestas confusas y obligó a separar acciones específicas para cada caso. Una segunda lección es que el clasificador NLU debe re-entrenarse con ejemplos extraídos del tráfico real una vez disponible, dado que las formulaciones espontáneas de los usuarios reales cubren un espacio léxico que el conjunto de \textit{stories} sintéticas no anticipa.

En el plano metodológico, la introducción de la suite de pruebas en una etapa temprana del desarrollo ha resultado decisiva. La primera versión de la suite contaba con 51 casos sobre la épica de asignaturas y arrojaba un 70,6\,\% de éxito; el sistema actual evalúa 159 casos con un 94,3\,\% de éxito. Sin la suite, las regresiones introducidas por cada cambio habrían sido muy difíciles de detectar a ojo. La lección complementaria es que cada decisión de diseño no trivial conviene registrarla en el momento, no al final: cuando se trata de reconstruir el porqué de un comportamiento al cabo de un mes, la memoria del autor ya no basta.

En el plano personal, gestionar un proyecto individual de la magnitud de un TFG exige imponerse una cadencia regular de trabajo, especialmente durante los meses en los que coexiste con asignaturas. La adopción de \textit{sprints} de tres o cuatro semanas ha funcionado bien como mecanismo de auto-disciplina, y el cierre de cada \textit{sprint} con un incremento desplegable ha sostenido la motivación a lo largo del curso.

\section{Análisis de desviación}\label{sec:conc-desviacion}

La planificación inicial recogida en el Capítulo~\ref{cap:planificacion} estimaba el esfuerzo total del proyecto en torno a 300 horas, repartidas entre tres fases (Inicio, Desarrollo y Cierre) y un conjunto de paquetes de trabajo. La ejecución real ha respetado el orden y la estructura previstos, pero ha presentado desviaciones puntuales que conviene comentar.

La fase de \textit{Inicio} ha consumido aproximadamente el esfuerzo previsto. La investigación tecnológica resultó algo más extensa de lo planeado por el tiempo dedicado a comparar modelos de lenguaje y a evaluar alternativas a Rasa, pero el diseño arquitectónico se completó dentro del margen estimado.

La fase de \textit{Desarrollo} ha sido la más larga y la que más ha sufrido desviaciones positivas, en el sentido de que se ha invertido más esfuerzo del estimado en algunos paquetes. La épica de profesorado, en concreto, requirió un ciclo adicional para resolver la combinatoria entre profesor, asignatura y grupo y para ajustar el atajo de coordinador y suplente, ambos bloqueos identificados durante el testing de aceptación. El \textit{pipeline} RAG también consumió más tiempo del previsto, principalmente por la fase de afinado de los \textit{prompts} de respuesta, que pasó por varias iteraciones hasta alcanzar la cifra del 91,9\,\% de profundidad.

La fase de \textit{Cierre}, por contra, se ha mantenido dentro del margen previsto. El piloto con usuarios reales y la auditoría manual de las 59 sesiones se ejecutaron en la ventana planificada (1 al 25 de abril de 2026), y la redacción de la memoria comenzó tan pronto como las épicas estuvieron estabilizadas, lo que ha evitado una concentración excesiva del esfuerzo en las últimas semanas.

Como recomendación de la guía de referencia~\cite{segura2024guiatfg}, las desviaciones se han mantenido siempre en sentido positivo: en ningún momento se ha invertido un volumen de horas inferior al estimado, lo que descartaría un cumplimiento meramente nominal del proyecto. La lección aprendida en este apartado es que la estimación inicial subestimaba sistemáticamente el coste de las tareas con dependencia de modelo de lenguaje (afinado de \textit{prompts}, evaluación manual de respuestas), debido a la naturaleza no determinista de ese tipo de tareas.

\section{Desarrollos futuros}\label{sec:conc-futuros}

El estado actual de Linceus Assistant corresponde a una fase 1, equivalente a un prototipo plenamente funcional pero con un alcance acotado a tres épicas y a la titulación GII-IS. El trabajo identificado para iteraciones posteriores se organiza en tres bloques: optimización del sistema actual, ampliación funcional y validación reforzada.

\subsection{Optimización del sistema actual}\label{subsec:conc-fut-optimizacion}

A lo largo del Capítulo~\ref{cap:pruebas} se han identificado cinco palancas técnicas con potencial para reducir la latencia, abaratar el coste y mejorar la robustez del sistema sin alterar su funcionalidad esencial. La primera es la \textbf{elección de modelo Gemini en función del tipo de llamada}: las tareas de \textit{Text-to-SQL} y de clasificación binaria pueden resolverse con modelos más pequeños y rápidos (familia Gemini Flash Lite o Gemma~3 de 4B parámetros), reservando modelos mayores para el renderizado natural de la respuesta. La estimación es que esta diferenciación reduciría a la mitad el coste anual y la cola larga de latencia.

La segunda palanca es la \textbf{sustitución selectiva de heurísticas léxicas por clasificadores ligeros basados en LLM o por similitud semántica con \textit{embeddings}}. La política de delegación descrita en la Sección~\ref{sec:politica-llm} optó deliberadamente por resolver con regex y diccionarios buena parte de las decisiones de menor entidad (detección de grupo, identificación de rol coordinador o suplente, regex de continuación elíptica como \textit{``los demás grupos''}, pronombres demostrativos para invertir prioridades de seguimiento entre profesor y asignatura, palabras clave para activar la rama RAG en consultas específicas) para mantener el coste y la latencia bajos y minimizar el riesgo de alucinaciones. Esta decisión es correcta para el alcance actual del prototipo, pero deja en evidencia dos limitaciones cuando se piensa en una versión productiva: las heurísticas son frágiles ante fraseos atípicos del estudiantado y están todas codificadas en español, lo que dificulta la internacionalización futura. La alternativa natural es una migración gradual a \textbf{clasificadores LLM ligeros} (modelos de la familia Gemini Flash Lite, una llamada por turno, temperatura cero) en los puntos donde la heurística falla con más frecuencia, o a \textbf{similitud semántica} con \textit{embeddings} multilingües cuando el problema es comparar la pregunta contra un conjunto pequeño de anclas conceptuales (por ejemplo, las palabras clave que activan RAG). Cualquiera de las dos vías convertiría varios de los casos actualmente excluidos como trabajo futuro en casos PASS, a coste de añadir entre 800~ms y 2~s por turno y de aceptar un margen no nulo de alucinación que hoy se evita por construcción.

La tercera palanca es una \textbf{caché de respuestas a las consultas más frecuentes}. La distribución del tráfico recogida en la tabla \texttt{conversation\_log} muestra que un puñado de preguntas (``¿qué es ADDA?'', ``horario de IS 2 grupo 1'', ``email de Galindo'') concentran un porcentaje significativo de los turnos. Una caché por \textit{hash} de pregunta y titulación reduciría a la mitad el coste anual y bajaría la latencia mediana al rango de uno o dos segundos.

La cuarta palanca es la \textbf{instrumentación de métricas en producción}: latencia por percentil, tasa de \textit{fallback}, tasa de \textit{out\_of\_scope} y \textit{feedback} explícito por turno. La tabla de \textit{feedback} ya está soportada parcialmente en la base de datos pero no se está explotando aún. Esta instrumentación permitiría detectar regresiones en producción mucho antes de que un usuario las reporte.

La quinta palanca es la \textbf{centralización de la lógica de seguimiento conversacional} en una acción genérica única. En la versión actual, cada épica implementa su propia heurística de herencia de \textit{slots} cuando detecta entidades vacías. La acción \texttt{ActionConsultaProfesor} hereda \texttt{ultimo\_profesor\_consultado} o \texttt{ultimo\_nombre\_asignatura} con una regla de prioridad invertible según pronombres demostrativos. La acción \texttt{ActionConsultaHorario} hereda \texttt{ultimo\_curso\_consultado} y \texttt{ultimo\_grupo\_consultado}, y cuando solo hay asignatura reciente delega en \texttt{ActionConsultaHorarioAsignatura}. Por último, \texttt{ActionSmartFallback} aplica un atajo regex específico para continuaciones de horario. Esta lógica está duplicada en tres ficheros con variantes sutiles, no cubre todos los seguimientos cruzados imaginables y deja casos de prueba como P-S01 y H-S01 (Capítulo~\ref{cap:pruebas}) en zona gris. La propuesta consiste en introducir un \texttt{ActionSeguimientoGenerico} que se invoque cuando NLU clasifique con baja confianza y la pregunta encaje con patrones de elipsis (regex barato sobre pronombres y conectores de continuación), consulte el \textit{slot} \texttt{ultima\_action\_ejecutada} junto con los \textit{slots} de contexto, y resuelva el destino mediante una \textbf{tabla de transiciones declarativa} en lugar de heurísticas dispersas. Solo si la tabla no resuelve un caso, escalaría al clasificador LLM actual del fallback. Este diseño elimina la duplicación, hace el sistema extensible (añadir una épica no requiere tocar las heurísticas de las anteriores), reduce el acoplamiento entre épicas (hoy materializado en \textit{imports} cruzados como el de \texttt{HorarioAsignatura}), y mantiene el coste cercano a cero en el caso común.

\subsection{Ampliación funcional}\label{subsec:conc-fut-ampliacion}

El alcance funcional del prototipo puede extenderse en varias direcciones. La más inmediata es la incorporación de las titulaciones del centro distintas a las tres de Ingeniería Informática actualmente cargadas, lo cual requiere principalmente trabajo de carga de datos en el panel de administración. A medio plazo, podría considerarse la incorporación de nuevos dominios al bot, como información sobre tutorías estructuradas, normativa académica, convocatorias de movilidad o \textbf{localización de espacios universitarios} (aulas, despachos, laboratorios), funcionalidad contemplada en el anteproyecto y aplazada en la Sección~\ref{sec:req-evolucion} por la falta de una fuente de datos estructurada y mantenible; su viabilidad dependerá de que la US publique planos o APIs de aforo en un formato procesable. Por último, la ampliación natural a otros centros de la Universidad de Sevilla pasa por validar el grado en que los \textit{prompts} y los flujos diseñados son transferibles cuando cambian los datos de origen.

\subsection{Validación reforzada}\label{subsec:conc-fut-validacion}

La validación realizada en la fase 1 cubre cinco capas funcionales y dos métricas no funcionales, pero deja fuera tres aspectos que conviene atender en una segunda iteración. El primero es la medición de la experiencia subjetiva del usuario mediante una escala estandarizada como SUS, que permitiría situar el sistema frente a otras herramientas universitarias. El segundo es la garantía de disponibilidad sostenida del servicio mediante un acuerdo de nivel de servicio y la correspondiente monitorización de uptime. El tercero es la implementación de la suite automática de pruebas para los \textit{endpoints} del panel de administración, cuyo diseño ya consta en el plan de aceptación pero cuyo desarrollo se ha pospuesto.

\section{Contribución al estado del arte}\label{sec:conc-contribucion-arte}

A la luz de los resultados obtenidos y de la revisión sistemática de chatbots educativos recopilada por Wollny et al.~\cite{wollny2021survey}, Linceus Assistant se sitúa en la intersección de dos líneas de investigación activas: los asistentes conversacionales para educación superior y los sistemas RAG sobre documentos de dominio cerrado. Frente a los cuatro sistemas universitarios comparables que han servido de referencia a lo largo de la memoria (AdmitBot~\cite{gupta2019admitbot}, EDUBOT~\cite{colace2018edubot}, Ranoliya~\cite{ranoliya2017chatbot} y Dibya--Sahoo~\cite{dibya2021chatbot}), el sistema se diferencia en tres aspectos concretos que dan forma a su aportación.

\paragraph{Ajuste iterativo y documentado en lugar de umbrales por convención.} Los porcentajes de éxito de los cuatro sistemas comparables se publican sin una descripción del procedimiento por el que se fijaron los umbrales internos del sistema (similitud semántica, fuzzy matching, fallback NLU). En Linceus, los tres umbrales críticos se acompañan de la observación cualitativa que los justifica: los valores 0,60 / 0,30 del \textit{fuzzy matching} de profesorado (Sección~\ref{sec:impl-shared}) responden al equilibrio entre tolerancia a typos moderados y rechazo de nombres cortos sin especificidad; la elevación del umbral de cambio de titulación a 0,85 (Sección~\ref{sec:impl-contexto}) corrige cambios silenciosos detectados en la revisión R3; y la configuración del par \texttt{threshold = 0,8} / \texttt{ambiguity\_threshold = 0,10} del clasificador de \textit{fallback} (Sección~\ref{subsec:dd-pipeline}) atiende a competencias entre intents observadas durante el desarrollo. Esta trazabilidad documental convierte cada cifra de la suite de pruebas en una decisión reproducible.

\paragraph{Arquitectura híbrida NLU + RAG con dos vías de resolución por dominio.} La revisión de Wollny~\cite{wollny2021survey} señala que la mayoría de chatbots educativos opta por un enfoque puramente conversacional sobre reglas o por una respuesta basada en LLM sin recuperación verificable. Linceus combina dentro de una misma \textit{custom action} una vía estructurada (\textit{Text-to-SQL} sobre Supabase) y una vía semántica (RAG sobre planes docentes con \texttt{pgvector}), con tres puntos de decisión por épica que escogen la rama adecuada en función del contenido de la pregunta (Sección~\ref{sec:impl-profesores}). Esta hibridación permite atender simultáneamente consultas que requieren precisión factual (créditos, horarios) y consultas que requieren interpretación textual (criterios de evaluación, coordinación), sin sacrificar ninguna de las dos.

\paragraph{Cierre del ciclo de vida del corpus mediante un panel sin código.} El tercer aporte responde a una crítica recurrente en la literatura: los prototipos académicos llegan al despliegue pero no a la operación sostenida, porque la actualización de la base de conocimiento exige al desarrollador. El panel Flask descrito en la Sección~\ref{sec:impl-admin} traslada la operación al perfil de mantenimiento, con flujos auditables para sincronizar el catálogo desde Sevius, enriquecer el directorio de profesorado desde \texttt{us.es} y vectorizar planes docentes con su pipeline completo de \textit{chunking}, \textit{embeddings} y \textit{rate limiting}. La suite \textit{pytest} de 24 casos sobre el panel (Sección~\ref{sec:pruebas-admin}) prueba que esa capa es operativa por sí misma, no un complemento decorativo.

Estas tres contribuciones no se publican como aportación al campo en un sentido fuerte, propio de una tesis doctoral, pero sí establecen un patrón concreto y reproducible para el desarrollo de asistentes universitarios. La distancia respecto a los prototipos publicados es prudente: la metodología de evaluación, el corpus utilizado y la naturaleza del dominio difieren entre proyectos y la comparación cuantitativa estricta no resulta procedente. Lo que sí cabe afirmar, a la vista de los datos del Capítulo~\ref{cap:pruebas}, es que Linceus alcanza tasas de éxito en los tres aspectos por encima de las cifras reportadas para el conjunto de los sistemas comparables.

\section{Consideraciones éticas}\label{sec:conc-etica}

El desarrollo de un asistente conversacional dirigido a estudiantes plantea responsabilidades que conviene explicitar más allá del simple cumplimiento legal. En el caso de Linceus, cuatro aspectos merecen comentario.

\paragraph{Veracidad de la información y riesgo de inducir decisiones erróneas.} El sistema puede influir, aunque sea de forma indirecta, en decisiones académicas relevantes para el alumno (elección de optativas, planificación de matrícula, identificación del docente coordinador para un trámite). Tres salvaguardas reducen este riesgo. La primera es que las respuestas sobre planes docentes están ancladas a documentos oficiales mediante el módulo RAG, no generadas a partir del conocimiento paramétrico del modelo. La segunda es el flujo de \textit{fallback} inteligente (Sección~\ref{sec:impl-fallback}), que declara explícitamente cuándo el sistema no encuentra información en lugar de inventar contenido. La tercera es la posibilidad de auditar a posteriori todas las conversaciones desde el panel de administración (RF-AD6), lo que permite identificar respuestas incorrectas y corregir el corpus o los \textit{prompts} responsables.

\paragraph{Privacidad y protección de datos.} El sistema se ha diseñado bajo el principio de minimización: la sesión es anónima (UUID generado del lado del cliente), no se solicita ninguna identificación personal del alumno y los datos académicos consultados (planes docentes, horarios y directorio público de profesorado) son de carácter público y proceden de fuentes oficiales (Sevius, \texttt{us.es}). La tabla \texttt{conversation\_log} almacena el contenido de las conversaciones para fines de auditoría y mejora, pero al no contener identificadores personales no constituye tratamiento de datos en el sentido del RGPD. El consentimiento del usuario para esta auditoría está cubierto por el aviso público del sitio donde se integra el \textit{widget}. La cuestión es relevante porque cualquier ampliación que introdujera datos identificables (autenticación universitaria, integración con expediente del alumno) exigiría una revisión completa de la base legal y la incorporación de mecanismos de consentimiento explícito.

\paragraph{Sesgo en los datos de entrenamiento.} El corpus NLU ha sido elaborado manualmente por el autor del trabajo, lo que introduce inevitablemente sesgos en la forma de formular las consultas: vocabulario, registros y patrones léxicos propios. La validación con tráfico real (Sección~\ref{sec:pruebas-trafico}) ha permitido detectar y completar formulaciones que el conjunto sintético no había anticipado, y el reentreno posterior con 117 ejemplos extraídos del piloto (Sección~\ref{sec:pruebas-nlu}) mitiga este sesgo en la versión actual. La diversificación sistemática del corpus mediante técnicas de \textit{data augmentation} (paráfrasis con LLM, inclusión de variantes regionales del español) figura entre las líneas de mejora.

\paragraph{Dependencia tecnológica y degradación responsable.} La fase generativa del sistema depende de la API de Gemini, propiedad de un tercero (Google). En caso de interrupción del servicio o de cambios en los términos de uso, el sistema degrada su comportamiento de forma controlada: el módulo RAG dispone de \textit{fallback} a búsqueda por palabras clave (Sección~\ref{subsec:impl-rag-busqueda}), las consultas estructurales más frecuentes siguen funcionando porque dependen exclusivamente de la base de datos relacional, y el cliente alternativo de Ollama (Sección~\ref{subsec:dd-gemini-ollama}) permite reconfigurar el sistema para una explotación local con un cambio de \textit{import}, sin reescribir la lógica de las acciones. Esta planificación de la degradación responde a la consideración ética de no dejar al alumno sin un servicio que pueda haber asumido como permanente.

\section{Reflexión final}\label{sec:conc-reflexion}

Mirar atrás al proyecto completo deja una impresión clara: la mayor dificultad de un sistema conversacional basado en modelos de lenguaje no reside en hacerlo funcionar, sino en hacerlo funcionar de forma fiable y trazable. Construir una primera versión que respondiera correctamente a un puñado de consultas demostrativas resultó relativamente rápido; alcanzar el punto en que cinco capas independientes de validación superan su umbral académico de referencia ha exigido un trabajo metódico de evaluación, corrección de defectos y registro de decisiones que ha constituido la mitad del esfuerzo del proyecto.

Esta experiencia sugiere que la incorporación de modelos de lenguaje a aplicaciones reales no se resuelve únicamente con buenas elecciones tecnológicas. Exige adoptar prácticas de ingeniería del software clásicas, en particular pruebas de aceptación, registro de decisiones y validación con usuarios, adaptadas al carácter no determinista del componente generativo. Linceus Assistant, en su versión actual, es un primer paso en esa dirección dentro del contexto universitario; las cinco palancas y las tres ampliaciones identificadas para la fase 2 deberían permitir convertirlo en una herramienta utilizable de forma sostenida por la comunidad de la ETSII.
